% ------------------------------------------------------------------------
% file `c3-exercise-1-solution-body.tex'
%   in folder `xsim/'
%
%     solution of type `exercise' with id `1'
%
% generated by the `solution' environment of the
%   `xsim' package v0.21 (2022/02/12)
% from source `c3' on 2023/06/08 on line 77
% ------------------------------------------------------------------------
    Évaluation :
    \begin{itemize}
        \item Trouver l'équation \hfill 1
        \item Factoriser \hfill 1
        \item Résoudre \hfill 1
        \item Interpréter \hfill 1
    \end{itemize}
    \begin{enumerate}
        \item Comme le triangle $ABC$ est rectangle en $A$, le théorème de Pythagore me permet d'écrire,
        \begin{align*}
            (x+1)^2 &= (x-1)^2 +  \left(\dfrac{x}{5}\right)^2\\
            x^2+2x+1 &= x^2 - 2x + 1 + \dfrac{x^2}{25}\\
            -4x + \dfrac{x^2}{25} &=0\\
            x^2 -100x &= 0
        \end{align*}
        \item Résolvons cette équation (mise en évidence et règle du produit nul)
        \begin{align*}
            x^2 - 100x &= 0\\
            x(x-100) &= 0\\
            x = 0 &\text{ ou }x=100
        \end{align*}
        \item Dans ce contexte, la solution $x=0$ est a rejetée car la longueur de côté d'un triangle ne peut être nul ou négative. Par conséquent, les longueurs des côtés de ce triangle sont \numlist{20;99;101}.
    \end{enumerate}

